% Generated from validated Special post-draft synthesis. Do not hand-edit.
\section*{Retrospective Signals}
\addcontentsline{toc}{section}{Retrospective Signals}
\smallskip
\noindent \textbf{フロンティア競争の軸がインタラクションへ広がった} 長い入力や画像を含む入力、提供形態までが同時に設計対象となり、単一の性能指標では捉えにくい競争へ移った。 \hfill{\footnotesize p.~\pageref{pkg:frontier-interface}}
\par
\smallskip
\noindent \textbf{公開モデルは単発の配布物から生態系へ変わった} 公開範囲、利用条件、専門化、規模、言語、効率の異なる系列が並走し、用途と計算資源に応じた選択肢が増えた。 \hfill{\footnotesize p.~\pageref{pkg:open-weight-ecosystem}}
\par
\smallskip
\noindent \textbf{生成の外側に実行と運用の層が立ち上がった} 外部機能を選ぶ仕組みと安全に実行する境界、限られた資源で適応・推論・配信する仕組みが、モデル能力とは別の設計問題として前面に出た。 \hfill{\footnotesize p.~\pageref{pkg:tool-action-layer} / p.~\pageref{pkg:runtime-adaptation}}
\par
\smallskip
\noindent \textbf{設計空間と生成メディアが同時に分岐した} 推論時探索、学習後最適化、検索制御、代替アーキテクチャが別々の研究面を広げる一方、画像・動画生成はモデル公開、製品統合、制作工程へと展開した。 \hfill{\footnotesize p.~\pageref{pkg:reasoning-retrieval-architecture} / p.~\pageref{pkg:generative-media}}
\par
\smallskip
\noindent \textbf{能力を測り、制御する仕組みが技術史の一部になった} 共通の比較基盤と、人間の選好を含む評価、能力向上と運用条件を結び付ける制御の枠組みが、強いモデルを扱う独立した技術層として現れた。 \hfill{\footnotesize p.~\pageref{pkg:evaluation-control} / p.~\pageref{pkg:annual-synthesis}}
\par
\medskip
\begin{claimboundary}[Retrospective scope]
本号は2023年を後日確認可能になった一次情報も用いて再構成するRetrospective Specialである。Coverage Periodと制作時点を同一視せず、vendor / project / author claimの境界は各記事で明示する。
\end{claimboundary}
\medskip
\tableofcontents
